\chapter{Base de Datos}

\section{Arquitectura General}

La base de datos utiliza \textbf{PostgreSQL 16} con \textbf{Drizzle ORM 0.41}
como capa de abstracci\'on. El esquema est\'a compuesto por 14 tablas
organizadas en torno al concepto de \textbf{multi-tenancy} por filas, donde
cada registro tiene un \texttt{clinic\_id} que lo asocia a una cl\'{\i}nica
espec\'{\i}fica.

\subsection{Diagrama de Entidades}

A continuaci\'on se presentan las entidades del sistema agrupadas por dominio:

\begin{itemize}[leftmargin=*]
    \item \textbf{Tenant root:} clinics
    \item \textbf{Staff:} users
    \item \textbf{Relaciones con clientes:} clients, patients
    \item \textbf{Atenci\'on m\'edica:} appointments, medical\_records, dental\_records
    \item \textbf{Vacunaci\'on:} vaccines, vaccination\_records
    \item \textbf{Facturaci\'on:} invoices, invoice\_items
    \item \textbf{Inventario:} products, suppliers
    \item \textbf{Comunicaci\'on:} notifications
\end{itemize}

\section{Definici\'on de Esquemas}

\subsection{clinics --- Tabla Ra\'{\i}z Multi-Tenant}

\begin{lstlisting}[style=typescript, caption=Esquema clinics]
export const clinics = pgTable('clinics', {
  id: serial('id').primaryKey(),
  name: text('name').notNull(),
  slug: text('slug').notNull().unique(),
  email: text('email'),
  phone: text('phone'),
  address: text('address'),
  settings: jsonb('settings').$type<{
    theme?: string;
    locale?: string;
    timezone?: string;
    invoicePrefix?: string;
    taxRate?: number;
  }>().default({}),
  active: boolean('active').notNull().default(true),
  createdAt: timestamp('created_at').notNull().defaultNow(),
  updatedAt: timestamp('updated_at').notNull().defaultNow(),
}, (table) => ({
  activeIdx: index('clinics_active_idx').on(table.active),
}));
\end{lstlisting}

\subsection{users --- Usuarios del Sistema}

\begin{lstlisting}[style=typescript, caption=Esquema users]
export const userRoleEnum = [
  'super_admin', 'admin', 'veterinarian',
  'technician', 'receptionist'
] as const;

export const users = pgTable('users', {
  id: serial('id').primaryKey(),
  clinicId: integer('clinic_id').notNull()
    .references(() => clinics.id),
  email: text('email').notNull().unique(),
  passwordHash: text('password_hash').notNull(),
  firstName: text('first_name').notNull(),
  lastName: text('last_name').notNull(),
  phone: text('phone'),
  role: text('role', { enum: userRoleEnum })
    .notNull().default('receptionist'),
  active: boolean('active').notNull().default(true),
  createdAt: timestamp('created_at').notNull().defaultNow(),
  updatedAt: timestamp('updated_at').notNull().defaultNow(),
}, (table) => ({
  clinicIdx: index('users_clinic_id_idx').on(table.clinicId),
  emailIdx: index('users_email_idx').on(table.email),
  roleIdx: index('users_role_idx').on(table.role),
}));
\end{lstlisting}

\subsection{clients --- Due\~nos de Mascotas}

\begin{lstlisting}[style=typescript, caption=Esquema clients]
export const clients = pgTable('clients', {
  clinicId: integer('clinic_id').notNull()
    .references(() => clinics.id),
  id: serial('id').primaryKey(),
  name: text('name').notNull(),
  email: text('email'),
  phone: text('phone').notNull(),
  address: text('address'),
  emergencyContact: text('emergency_contact'),
  notes: text('notes'),
  metadata: jsonb('metadata'),
  active: boolean('active').notNull().default(true),
  createdAt: timestamp('created_at').notNull().defaultNow(),
  updatedAt: timestamp('updated_at').notNull().defaultNow(),
}, (table) => ({
  clinicIdx:  index('clients_clinic_id_idx').on(table.clinicId),
  emailIdx:   index('clients_email_idx').on(table.email),
  phoneIdx:   index('clients_phone_idx').on(table.phone),
  nameIdx:    index('clients_name_idx').on(table.name),
}));
\end{lstlisting}

\subsection{patients --- Pacientes (Mascotas)}

\begin{lstlisting}[style=typescript, caption=Esquema patients]
export const speciesEnum = ['dog','cat','rabbit','bird','reptile','other'] as const;
export const genderEnum = ['male','female'] as const;

export const patients = pgTable('patients', {
  id: serial('id').primaryKey(),
  clinicId: integer('clinic_id').notNull()
    .references(() => clinics.id),
  clientId: integer('client_id').notNull()
    .references(() => clients.id),
  name: text('name').notNull(),
  species: text('species', { enum: speciesEnum }).notNull(),
  breed: text('breed'),
  gender: text('gender', { enum: genderEnum }).notNull(),
  birthDate: date('birth_date'),
  color: text('color'),
  weight: jsonb('weight')
    .$type<{ value: number; date: string }[]>().default([]),
  microchip: text('microchip'),
  photoUrl: text('photo_url'),
  allergies: text('allergies'),
  notes: text('notes'),
  active: boolean('active').notNull().default(true),
  createdAt: timestamp('created_at').notNull().defaultNow(),
  updatedAt: timestamp('updated_at').notNull().defaultNow(),
}, (table) => ({
  clinicIdx:  index('patients_clinic_id_idx').on(table.clinicId),
  clientIdx:  index('patients_client_id_idx').on(table.clientId),
  nameIdx:    index('patients_name_idx').on(table.name),
  speciesIdx: index('patients_species_idx').on(table.species),
}));
\end{lstlisting}

\subsection{appointments --- Citas}

\begin{lstlisting}[style=typescript, caption=Esquema appointments]
export const appointmentStatusEnum = [
  'scheduled', 'confirmed', 'checked_in', 'in_exam',
  'checked_out', 'cancelled', 'no_show',
] as const;

export const appointmentTypeEnum = [
  'consultation', 'vaccination', 'surgery', 'emergency',
  'grooming', 'boarding', 'follow_up', 'other',
] as const;

export const appointments = pgTable('appointments', {
  id: serial('id').primaryKey(),
  clinicId: integer('clinic_id').notNull()
    .references(() => clinics.id),
  patientId: integer('patient_id').notNull()
    .references(() => patients.id),
  clientId: integer('client_id').notNull()
    .references(() => clients.id),
  vetId: integer('vet_id').references(() => users.id),
  date: date('date').notNull(),
  startTime: time('start_time').notNull(),
  endTime: time('end_time').notNull(),
  type: text('type', { enum: appointmentTypeEnum })
    .notNull().default('consultation'),
  status: text('status', { enum: appointmentStatusEnum })
    .notNull().default('scheduled'),
  notes: text('notes'),
  room: text('room'),
  createdBy: integer('created_by').references(() => users.id),
  createdAt: timestamp('created_at').notNull().defaultNow(),
  updatedAt: timestamp('updated_at').notNull().defaultNow(),
}, (table) => ({
  clinicIdx:    index('appointments_clinic_id_idx')
                .on(table.clinicId),
  patientIdx:   index('appointments_patient_id_idx')
                .on(table.patientId),
  clientIdx:    index('appointments_client_id_idx')
                .on(table.clientId),
  vetIdx:       index('appointments_vet_id_idx')
                .on(table.vetId),
  dateIdx:      index('appointments_date_idx')
                .on(table.date),
  statusIdx:    index('appointments_status_idx')
                .on(table.status),
  dateStatusIdx: index('appointments_date_status_idx')
                .on(table.date, table.status),
  startTimeIdx: index('appointments_start_time_idx')
                .on(table.startTime),
}));
\end{lstlisting}

\subsection{medical\_records --- Historial Cl\'{\i}nico (SOAP)}

\begin{lstlisting}[style=typescript, caption=Esquema medical\_records]
export const medicalRecords = pgTable('medical_records', {
  id: serial('id').primaryKey(),
  clinicId: integer('clinic_id').notNull()
    .references(() => clinics.id),
  patientId: integer('patient_id').notNull()
    .references(() => patients.id),
  appointmentId: integer('appointment_id')
    .references(() => appointments.id),
  subjective: text('subjective'),
  objective: text('objective'),
  assessment: text('assessment'),
  plan: text('plan'),
  vitals: jsonb('vitals').$type<{
    weight?: number;
    temperature?: number;
    heartRate?: number;
    respiratoryRate?: number;
  }>().default({}),
  diagnosis: text('diagnosis'),
  createdBy: integer('created_by')
    .references(() => users.id),
  createdAt: timestamp('created_at').notNull().defaultNow(),
  updatedAt: timestamp('updated_at').notNull().defaultNow(),
}, (table) => ({
  clinicIdx:    index('medical_records_clinic_id_idx')
                .on(table.clinicId),
  patientIdx:   index('medical_records_patient_id_idx')
                .on(table.patientId),
  appointmentIdx: index('medical_records_appointment_id_idx')
                  .on(table.appointmentId),
  createdAtIdx: index('medical_records_created_at_idx')
                .on(table.createdAt),
}));
\end{lstlisting}

\subsection{dental\_records --- Odontograma}

\begin{lstlisting}[style=typescript, caption=Esquema dental\_records]
export type ToothStatusValue =
  | 'healthy' | 'fractured' | 'missing' | 'worn'
  | 'calculus' | 'gingivitis' | 'mobility'
  | 'retained_deciduous' | 'root_remnant' | 'caries'
  | 'pulp_exposure' | 'other';

export interface ToothCondition {
  toothNumber: number;
  status: ToothStatusValue;
  notes?: string;
}

export const dentalRecords = pgTable('dental_records', {
  id: serial('id').primaryKey(),
  clinicId: integer('clinic_id').notNull()
    .references(() => clinics.id),
  patientId: integer('patient_id').notNull()
    .references(() => patients.id),
  teeth: jsonb('teeth')
    .$type<ToothCondition[]>().default([]),
  notes: jsonb('notes')
    .$type<Record<string, string>>().default({}),
  createdBy: integer('created_by')
    .references(() => users.id),
  createdAt: timestamp('created_at').notNull().defaultNow(),
  updatedAt: timestamp('updated_at').notNull().defaultNow(),
}, (table) => ({
  clinicIdx:  index('dental_records_clinic_id_idx')
              .on(table.clinicId),
  patientIdx: index('dental_records_patient_id_idx')
              .on(table.patientId),
}));
\end{lstlisting}

\subsection{vaccines --- Cat\'alogo de Vacunas}

\begin{lstlisting}[style=typescript, caption=Esquema vaccines]
export const vaccines = pgTable('vaccines', {
  id: serial('id').primaryKey(),
  clinicId: integer('clinic_id').notNull()
    .references(() => clinics.id),
  name: text('name').notNull(),
  description: text('description'),
  species: text('species').notNull(),
  frequencyDays: integer('frequency_days')
    .notNull().default(365),
  isCore: boolean('is_core').notNull().default(true),
  active: boolean('active').notNull().default(true),
}, (table) => ({
  clinicIdx:  index('vaccines_clinic_id_idx')
              .on(table.clinicId),
  speciesIdx: index('vaccines_species_idx')
              .on(table.species),
}));
\end{lstlisting}

\subsection{vaccination\_records --- Registros de Vacunaci\'on}

\begin{lstlisting}[style=typescript, caption=Esquema vaccination\_records]
export const vaccinationRecords = pgTable('vaccination_records', {
  id: serial('id').primaryKey(),
  clinicId: integer('clinic_id').notNull()
    .references(() => clinics.id),
  patientId: integer('patient_id').notNull()
    .references(() => patients.id),
  vaccineId: integer('vaccine_id').notNull()
    .references(() => vaccines.id),
  appliedDate: date('applied_date').notNull(),
  nextDueDate: date('next_due_date'),
  batchNumber: text('batch_number'),
  applicationSite: text('application_site'),
  doseNumber: integer('dose_number').default(1),
  administeredBy: integer('administered_by')
    .references(() => users.id),
  notes: text('notes'),
  createdAt: timestamp('created_at').notNull().defaultNow(),
  updatedAt: timestamp('updated_at').notNull().defaultNow(),
}, (table) => ({
  clinicIdx:    index('vaccination_records_clinic_id_idx')
                .on(table.clinicId),
  patientIdx:   index('vaccination_records_patient_id_idx')
                .on(table.patientId),
  vaccineIdx:   index('vaccination_records_vaccine_id_idx')
                .on(table.vaccineId),
  patientDueIdx: index('vaccination_records_patient_due_idx')
                 .on(table.patientId, table.nextDueDate),
}));
\end{lstlisting}

\subsection{invoices e invoice\_items --- Facturaci\'on}

\begin{lstlisting}[style=typescript, caption=Esquemas de facturacion]
export const invoiceStatusEnum = [
  'draft', 'issued', 'paid', 'cancelled'
] as const;

export const invoices = pgTable('invoices', {
  id: serial('id').primaryKey(),
  clinicId: integer('clinic_id').notNull()
    .references(() => clinics.id),
  clientId: integer('client_id').notNull()
    .references(() => clients.id),
  appointmentId: integer('appointment_id')
    .references(() => appointments.id),
  invoiceNumber: text('invoice_number').notNull().unique(),
  subtotal: decimal('subtotal', { precision: 10, scale: 2 })
    .notNull().default('0'),
  tax: decimal('tax', { precision: 10, scale: 2 })
    .notNull().default('0'),
  total: decimal('total', { precision: 10, scale: 2 })
    .notNull().default('0'),
  status: text('status', { enum: invoiceStatusEnum })
    .notNull().default('draft'),
  notes: text('notes'),
  createdAt: timestamp('created_at').notNull().defaultNow(),
  updatedAt: timestamp('updated_at').notNull().defaultNow(),
}, (table) => ({
  clinicIdx:      index('invoices_clinic_id_idx')
                  .on(table.clinicId),
  clientIdx:      index('invoices_client_id_idx')
                  .on(table.clientId),
  statusIdx:      index('invoices_status_idx')
                  .on(table.status),
  createdAtIdx:   index('invoices_created_at_idx')
                  .on(table.createdAt),
  clinicStatusIdx: index('invoices_clinic_status_idx')
                   .on(table.clinicId, table.status),
}));

export const invoiceItems = pgTable('invoice_items', {
  id: serial('id').primaryKey(),
  invoiceId: integer('invoice_id').notNull()
    .references(() => invoices.id),
  description: text('description').notNull(),
  quantity: integer('quantity').notNull().default(1),
  unitPrice: decimal('unit_price', { precision: 10, scale: 2 })
    .notNull(),
  total: decimal('total', { precision: 10, scale: 2 })
    .notNull(),
});
\end{lstlisting}

\subsection{products --- Inventario}

\begin{lstlisting}[style=typescript, caption=Esquema products]
export const productTypeEnum = [
  'medication', 'supply', 'food', 'service'
] as const;

export const products = pgTable('products', {
  id: serial('id').primaryKey(),
  clinicId: integer('clinic_id').notNull()
    .references(() => clinics.id),
  name: text('name').notNull(),
  description: text('description'),
  type: text('type', { enum: productTypeEnum }).notNull(),
  sku: text('sku').unique(),
  price: decimal('price', { precision: 10, scale: 2 }).notNull(),
  costPrice: decimal('cost_price', { precision: 10, scale: 2 }),
  stockQuantity: integer('stock_quantity').notNull().default(0),
  reorderPoint: integer('reorder_point').default(10),
  active: boolean('active').notNull().default(true),
  createdAt: timestamp('created_at').notNull().defaultNow(),
  updatedAt: timestamp('updated_at').notNull().defaultNow(),
}, (table) => ({
  clinicIdx: index('products_clinic_id_idx').on(table.clinicId),
  nameIdx:   index('products_name_idx').on(table.name),
  typeIdx:   index('products_type_idx').on(table.type),
}));
\end{lstlisting}

\subsection{suppliers --- Proveedores}

\begin{lstlisting}[style=typescript, caption=Esquema suppliers]
export const suppliers = pgTable('suppliers', {
  clinicId: integer('clinic_id').notNull()
    .references(() => clinics.id),
  id: serial('id').primaryKey(),
  name: text('name').notNull(),
  contact: text('contact'),
  email: text('email'),
  phone: text('phone'),
  address: text('address'),
  active: boolean('active').notNull().default(true),
  createdAt: timestamp('created_at').notNull().defaultNow(),
  updatedAt: timestamp('updated_at').notNull().defaultNow(),
}, (table) => ({
  clinicIdx: index('suppliers_clinic_id_idx').on(table.clinicId),
  nameIdx:   index('suppliers_name_idx').on(table.name),
}));
\end{lstlisting}

\subsection{notifications --- Notificaciones}

\begin{lstlisting}[style=typescript, caption=Esquema notifications]
export const notifications = pgTable('notifications', {
  id: serial('id').primaryKey(),
  userId: integer('user_id').notNull()
    .references(() => users.id),
  title: text('title').notNull(),
  message: text('message').notNull(),
  type: text('type', {
    enum: ['appointment','invoice','client','patient','system']
  }).notNull().default('system'),
  link: text('link'),
  read: boolean('read').notNull().default(false),
  createdAt: timestamp('created_at').notNull().defaultNow(),
}, (table) => ({
  userIdx:    index('notifications_user_id_idx')
              .on(table.userId),
  readIdx:    index('notifications_read_idx')
              .on(table.read),
  createdAtIdx: index('notifications_created_at_idx')
                .on(table.createdAt),
  userReadIdx: index('notifications_user_read_idx')
               .on(table.userId, table.read),
}));
\end{lstlisting}

\section{Relaciones entre Tablas}

Drizzle ORM permite definir relaciones expl\'{\i}citas que habilitan
consultas con m\'ultiples joins de forma tipo-safe:

\begin{lstlisting}[style=typescript, caption=Relaciones definidas con drizzle-orm/relations]
// Pacientes -> Cliente (cada mascota pertenece a un dueno)
export const patientsRelations = relations(patients, ({ one, many }) => ({
  client: one(clients, {
    fields: [patients.clientId],
    references: [clients.id],
  }),
  vaccinationRecords: many(vaccinationRecords),
}));

// Citas -> Paciente, Cliente, Veterinario
export const appointmentsRelations = relations(appointments, ({ one }) => ({
  patient: one(patients, {
    fields: [appointments.patientId],
    references: [patients.id],
  }),
  client: one(clients, {
    fields: [appointments.clientId],
    references: [clients.id],
  }),
  veterinarian: one(users, {
    fields: [appointments.vetId],
    references: [users.id],
  }),
}));

// Historial Clinico -> Paciente, Cita, Usuario (creador)
export const medicalRecordsRelations = relations(medicalRecords, ({ one }) => ({
  patient: one(patients, {
    fields: [medicalRecords.patientId],
    references: [patients.id],
  }),
  appointment: one(appointments, {
    fields: [medicalRecords.appointmentId],
    references: [appointments.id],
  }),
  createdByUser: one(users, {
    fields: [medicalRecords.createdBy],
    references: [users.id],
  }),
}));
\end{lstlisting}

\section{\'Indices}

La migraci\'on \texttt{0004\_brainy\_queen\_noir.sql} crea 38 \'{\i}ndices
distribuidos as\'{\i}:

\begin{longtable}{|l|l|l|}
\hline
\textbf{Tabla} & \textbf{\'Indices} & \textbf{Tipo} \\
\hline
\endhead
users & clinic\_id, email, role & B-tree \\
clients & clinic\_id, email, phone, name & B-tree \\
patients & clinic\_id, client\_id, name, species & B-tree \\
appointments & clinic\_id, patient\_id, client\_id, vet\_id, date,
               status, (date, status), start\_time & B-tree \\
medical\_records & clinic\_id, patient\_id, appointment\_id, created\_at & B-tree \\
products & clinic\_id, name, type & B-tree \\
suppliers & clinic\_id, name & B-tree \\
invoices & clinic\_id, client\_id, status, created\_at,
           (clinic\_id, status) & B-tree \\
notifications & user\_id, read, created\_at, (user\_id, read) & B-tree \\
clinics & active & B-tree \\
vaccines & clinic\_id, species & B-tree \\
vaccination\_records & clinic\_id, patient\_id, vaccine\_id,
                       (patient\_id, next\_due\_date) & B-tree \\
dental\_records & clinic\_id, patient\_id & B-tree \\
\hline
\end{longtable}

\section{Migraciones}

Las migraciones se gestionan con \textbf{Drizzle Kit} y se almacenan en
\texttt{packages/db/drizzle}. El orden de migraciones es:

\begin{lstlisting}[style=json, caption=\_journal.json - orden de migraciones]
{
  "entries": [
    { "idx": 0, "tag": "0000_jazzy_penance" },      // Schema inicial
    { "idx": 1, "tag": "0001_black_shard" },         // Facturacion
    { "idx": 2, "tag": "0002_notifications" },       // Notificaciones
    { "idx": 3, "tag": "0003_multi-tenant" },        // Multi-tenant
    { "idx": 4, "tag": "0004_brainy_queen_noir" },   // 38 indices
    { "idx": 5, "tag": "0005_square_songbird" },     // Vacunas
    { "idx": 6, "tag": "0006_young_dental_records" } // Odontograma
  ]
}
\end{lstlisting}

\subsection{Migraci\'on 0000 --- Esquema Inicial}

Crea las tablas: \texttt{users}, \texttt{clients}, \texttt{patients},
\texttt{appointments}, \texttt{medical\_records}, \texttt{products},
\texttt{suppliers}. Sin columna \texttt{clinic\_id} (pre-multi-tenant).

\subsection{Migraci\'on 0001 --- Facturaci\'on}

Agrega \texttt{invoices} e \texttt{invoice\_items}.

\subsection{Migraci\'on 0002 --- Notificaciones}

Agrega \texttt{notifications}.

\subsection{Migraci\'on 0003 --- Multi-Tenant}

Migraci\'on cr\'{\i}tica que:

\begin{enumerate}
    \item Crea la tabla \texttt{clinics} e inserta la cl\'{\i}nica por defecto.
    \item Agrega \texttt{clinic\_id} nullable a 8 tablas existentes.
    \item Actualiza registros existentes con el ID de la cl\'{\i}nica por defecto.
    \item Establece NOT NULL y agrega llaves for\'aneas.
\end{enumerate}

\subsection{Migraci\'on 0004 --- \'Indices}

Crea 38 \'{\i}ndices B-tree para optimizar consultas.

\subsection{Migraciones 0005 y 0006}

Agregan los m\'odulos de vacunaci\'on y odontograma respectivamente.

\section{Datos Semilla}

El archivo \texttt{packages/db/src/seed.ts} (534 l\'{\i}neas) genera datos
de prueba realistas:

\begin{itemize}[leftmargin=*]
    \item \textbf{1 cl\'{\i}nica:} ``Cl\'{\i}nica Veterinaria Central'' con
          configuraci\'on mexicana.
    \item \textbf{5 usuarios:} super\_admin, 2 veterinarios, recepcionista,
          t\'ecnico (pass: \texttt{admin123}).
    \item \textbf{15 clientes:} Due\~nos con datos completos.
    \item \textbf{22 pacientes:} 13 perros, 6 gatos, 2 aves, 2 conejos.
    \item \textbf{10 vacunas:} Con frecuencias y especies objetivo.
    \item \textbf{20 registros de vacunaci\'on.}
    \item \textbf{5 proveedores} y \textbf{30 productos.}
    \item \textbf{34 citas:} Con estados variados (programadas, confirmadas,
          canceladas, etc.).
    \item \textbf{8 registros m\'edicos} en formato SOAP.
    \item \textbf{2 registros dentales} con condiciones de dientes.
    \item \textbf{12 facturas} con partidas y estados variados.
    \item \textbf{11 notificaciones.}
\end{itemize}

El script limpia todas las tablas en orden de dependencias antes de insertar.
